\urlstyle{same}
\renewcommand\UrlFont{\normalfont}
\section{Технический проект}
\subsection{Архитектура модуля взаимодействия с пользователем}

Разрабатываемая программная часть представляет собой серверный модуль взаимодействия с пользователем для онлайн-площадки аренды вещей. Модуль реализован в виде набора Django-приложений и предоставляет API для работы с чатами, уведомлениями, договорами аренды и демонстрационным подписанием договоров электронной подписью.

Модуль является частью клиент-серверной архитектуры онлайн-площадки. Клиентская часть обращается к серверу по REST API для выполнения основных операций: получения списка диалогов, отправки сообщений, просмотра уведомлений, получения договора по бронированию и подписания договора. Для доставки событий в режиме реального времени используются WebSocket-соединения. Через них клиентская часть получает новые сообщения, события прочтения сообщений и уведомления пользователя.

Серверная часть реализована на языке Python с использованием фреймворка Django. Для построения REST API используется Django REST Framework, а для поддержки WebSocket-соединений применяется Django Channels. Хранение данных выполняется средствами Django ORM в реляционной базе данных. Для обмена событиями между WebSocket-подключениями в проекте предусмотрено использование Redis как промежуточного слоя сообщений.

В рамках модуля выделены три основные подсистемы:
\begin{itemize}
    \item подсистема чатов, обеспечивающая создание диалогов, хранение сообщений, отправку изображений и отметку сообщений как прочитанных;
    \item подсистема уведомлений, обеспечивающая создание системных уведомлений, хранение их статуса, доставку через WebSocket и отправку email-оповещений;
    \item подсистема договоров, обеспечивающая формирование PDF-договора аренды по данным бронирования и демонстрационное подписание договора арендатором и арендодателем.
\end{itemize}

Общая схема взаимодействия компонентов приведена на рисунке \ref{uml:image}

\begin{figure}[h]
\center{\includegraphics[width=0.7\linewidth]{canvas}}
\caption{Общая схема взаимодействия компонентов модуля взаимодействия с пользователем}
\label{uml:image}
\end{figure}

Выбор такой архитектуры обусловлен требованиями к разрабатываемому модулю. Серверная часть централизует бизнес-логику, проверку прав доступа, сохранение истории переписки, создание уведомлений и формирование договоров. Использование REST API позволяет отделить клиентский интерфейс от серверной логики, а применение WebSocket обеспечивает своевременную доставку сообщений и уведомлений без постоянного опроса сервера.


\subsection{Обоснование выбора технологии проектирования}
В настоящее время рынок программных решений в рассматриваемой области предоставляет множество инструментов, которые позволяют успешно реализовать серверную часть веб-приложения, в том числе модули для отправки уведомлений, организации чата между пользователями и формирования PDF-договоров.

\subsubsection{Описание используемых технологий и языков программирования}

При разработке серверной части проекта использовался язык программирования Python. На нём написана основная бизнес-логика системы, включая обработку бронирований, уведомлений, чатов и генерацию договоров аренды.

В качестве основного веб-фреймворка применялся Django. Он использовался для построения структуры проекта, описания моделей данных, работы с базой данных, настройки маршрутов и реализации серверной логики. Для создания REST API использовался Django REST Framework, позволяющий сериализовать данные, обрабатывать HTTP-запросы и настраивать права доступа к ресурсам.

Для реализации обмена сообщениями и уведомлений в реальном времени применялась технология WebSocket. На стороне серверной части она реализована с использованием Django Channels. Данная технология позволила организовать постоянное соединение между клиентом и сервером для доставки новых сообщений в чатах и системных уведомлений в режиме реального времени.

В качестве промежуточного слоя для WebSocket-сообщений использовался Redis. Он выполняет роль посредника сообщений между серверными процессами и обеспечивает доставку событий в соответствующие пользовательские каналы.

Для хранения данных использовалась реляционная система управления базами данных PostgreSQL. В ней хранятся сведения о пользователях, объявлениях, бронированиях, уведомлениях, чатах, сообщениях и договорах. 

Для формирования PDF-договоров применялась библиотека Pillow. С её помощью текст договора, сформированный на основе данных бронирования, преобразуется в PDF-файл и сохраняется в файловом хранилище проекта.

Для отправки email-уведомлений использовался стандартный почтовый механизм Django через SMTP. Он позволяет отправлять пользователям сообщения о новых заявках, подтверждении бронирования, оплате, новых сообщениях и напоминаниях о возврате предмета аренды.

Развёртывание проекта выполнялось с использованием Docker и Docker Compose. Эти инструменты позволяют запускать приложение, базу данных PostgreSQL, Redis и другие компоненты в изолированных контейнерах.

\subsubsection{Django как серверная платформа}
Django выбран как основной backend-фреймворк по нескольким причинам. Во-первых, он предоставляет модульную структуру приложений. В проекте это выражено в наличии отдельных приложений notifications, chats и contracts, каждое из которых отвечает за собственную предметную область. 

Во-вторых, Django ORM позволяет описывать модели данных декларативно и работать с ними через объектный интерфейс. Это снижает количество ручного SQL-кода и уменьшает риск ошибок при работе со связанными объектами.

В-третьих, Django хорошо подходит для проектов, где важна авторизация и контроль доступа. В системе аренды значительная часть данных является персональной: переписка, уведомления, сведения о сделках, договоры и подписи. Поэтому каждый API-метод должен учитывать текущего пользователя и проверять его роль в конкретной операции.

Ограничением Django является то, что классическая модель обработки HTTP-запросов не предназначена для постоянных двусторонних соединений. Для решения этой задачи используется Django Channels, расширяющий проект поддержкой ASGI и WebSocket.

\subsubsection{Django REST Framework как средство построения API}
Django REST Framework применяется для реализации REST API подсистем. Его использование позволяет привести интерфейсы разных модулей к единому виду. Например, уведомления выдаются через read-only viewset, договоры через viewset с дополнительными действиями, а чаты через APIView-классы, где требуется более гибкая логика создания диалогов и сообщений.

Основные преимущества DRF в проекте:
\begin{itemize}
\item сериализация моделей в JSON;
\item проверка входных данных;
\item настройка прав доступа;
\item единый формат HTTP-ответов;
\item удобное тестирование API через APIClient.
\end{itemize}

Для подсистемы договоров сериализатор дополнительно формирует вычисляемые поля: ссылку на файл договора, статус подписания, роль текущего пользователя и признак того, подписал ли текущий пользователь договор. Это позволяет клиентской части получать уже подготовленную структуру без необходимости самостоятельно анализировать все поля модели.

\subsubsection{Django Channels и Redis для событий реального времени}
Чат и уведомления требуют немедленной доставки событий. Например, если арендатор отправил сообщение арендодателю, второму пользователю необходимо получить его сразу. Аналогично при появлении нового уведомления система должна передать его адресату без задержки.
Для этого используется WebSocket. В проекте определены два основных типа соединений:
\begin{itemize}
\item /ws/chat/<chat-id>/ - соединение с конкретным чатом;
\item /ws/notifications/ - соединение с персональным каналом уведомлений пользователя.
\end{itemize}

При подключении пользователь проходит проверку авторизации. Для чата дополнительно проверяется, что пользователь является участником соответствующего диалога. Это предотвращает доступ к чужой переписке даже при известном идентификаторе чата.
Redis используется как channel layer. Сервер отправляет событие в группу, например chat-5 или notifications-12, а Channels доставляет его всем активным соединениям, подписанным на эту группу.

\subsubsection{PostgreSQL как СУБД}
PostgreSQL применяется как основное хранилище данных. Для разработанных подсистем важна ссылочная целостность, так как данные не существуют изолированно. Уведомление связано с пользователем, чат связан с участниками и при необходимости с предметом аренды, сообщение связано с чатом и отправителем, договор связан с бронированием.

Такая структура естественно выражается через внешние ключи и связи многие-ко-многим. PostgreSQL также обеспечивает транзакционность операций. Это важно, например, при подписании договора, когда необходимо обновить данные подписи, статус договора и затем пересформировать PDF-файл.

\subsubsection{Pillow для генерации PDF-договоров}
Для автоматического формирования договора аренды используется программная генерация PDF. Модуль договора собирает текст на основании данных бронирования: номера документа, предмета аренды, арендодателя, арендатора, дат, суммы и статуса. После этого текст переносится на страницы PDF и сохраняется в файловое поле модели.

Выбор Pillow объясняется простотой реализации и достаточностью для целей проекта. Основная задача состоит в автоматическом формировании читаемого документа, который отражает данные сделки и статусы подписания сторон.

\subsection{Описание разработанных модулей}
В логике проекта выделены следующие функциональные модули:
\begin{enumerate}
\item Модуль уведомлений. Отвечает за создание системных уведомлений, хранение их статуса, передачу уведомлений по WebSocket и отправку внешних сообщений по email.
\item Модуль чатов. Обеспечивает создание диалогов между арендатором и арендодателем, хранение сообщений, передачу новых сообщений в реальном времени, отметку сообщений как прочитанных и привязку переписки к конкретному предмету аренды.
\item Модуль договоров. Формирует договор аренды на основании данных бронирования, предмета аренды и участников сделки. Модуль создаёт PDF-файл, хранит его в файловом хранилище проекта и поддерживает демонстрационную процедуру подписания договора обеими сторонами.
\end{enumerate}


Такое разделение позволяет рассматривать разработанные части как относительно независимые сервисные подсистемы внутри единого Django-проекта. Несмотря на то что физически они развёрнуты в одном приложении, каждая подсистема имеет собственные модели, сериализаторы, маршруты, обработчики и тесты. Это снижает связность кода и упрощает дальнейшее выделение модулей в отдельные микросервисы при росте проекта.

\subsection{Структура разработанных модулей}
Модуль notifications включает:
\begin{itemize}
\item модель уведомления;
\item сервисные функции создания уведомлений;
\item функции отправки email;
\item WebSocket consumer для доставки уведомлений;
\item REST API для списка уведомлений и изменения статуса прочтения;
\item management-команду для отправки напоминаний о возврате.
\end{itemize}

Модуль chats включает:
\begin{itemize}
\item модель чата;
\item модель сообщения;
\item API создания и получения диалогов;
\item API отправки сообщений;
\item WebSocket consumer для обмена сообщениями в реальном времени;
\item логику отметки сообщений как прочитанных;
\item поддержку изображений в сообщениях.
\end{itemize}

Модуль contracts включает:
\begin{itemize}
\item модель договора;
\item генерацию номера документа;
\item формирование PDF-файла;
\item API получения договора по бронированию;
\item демонстрационное подписание договора арендатором и арендодателем;
\item проверку доступа к договору только для участников сделки.
\end{itemize}

Такое разделение позволяет поддерживать код по предметным областям. Логика уведомлений не смешивается с логикой чатов, а генерация договоров вынесена в отдельный модуль, связанный с бронированием через отдельную модель.

\subsection{Серверная часть}
Серверная часть проекта выполняет роль центрального слоя бизнес-логики. Она принимает запросы от клиента, проверяет права доступа, изменяет состояние базы данных и инициирует дополнительные действия. Например, создание бронирования приводит к созданию уведомления арендодателю, подтверждение бронирования создаёт уведомление арендатору, а отправка сообщения в чат создаёт уведомление получателю.

Для маршрутизации HTTP-запросов используются URL-конфигурации Django. Разработанные модули подключаются к общему API проекта по следующим префиксам:
\begin{itemize}
\item /api/notifications/ - уведомления;
\item /api/chats/ - диалоги, сообщения и справка;
\item /api/contracts/ - договоры аренды;
\item /ws/chat/<chat-id>/ - WebSocket конкретного чата;
\item /ws/notifications/ - WebSocket уведомлений пользователя.
\end{itemize}

Серверная часть также обеспечивает связь с существующими модулями проекта: пользователями, товарами и бронированиями. Это необходимо, так как уведомления, чаты и договоры возникают не сами по себе, а как результат действий в основном бизнес-процессе аренды.
\subsubsection{Маршруты API}
Все API-маршруты имеют общий префикс /api/. Авторизация выполняется по JWT-токену, передаваемому в заголовке Authorization: Bearer <token>. WebSocket-маршруты объявлены в asgi.py и используют JWT из строки запроса (?token=...), верифицируемый JWTAuthMiddleware. Основные HTTP-маршруты серверной части представлены в таблице 3.1.
\begin{xltabular}{\textwidth}{|
>{\raggedright\arraybackslash\hsize=0.85\hsize\linewidth=\hsize}X|
>{\centering\arraybackslash\hsize=0.35\hsize\linewidth=\hsize}X|
>{\raggedright\arraybackslash\hsize=0.70\hsize\linewidth=\hsize}X|
>{\raggedright\arraybackslash\hsize=1.10\hsize\linewidth=\hsize}X|}
 \caption{Основные HTTP-маршруты API \label{analiz:table}}\\ \hline
   \centrow Маршрут & \centrow Метод & \centrow Входные данные & \centrow Результат \\ \hline
 \endfirsthead

\continuecaption{Продолжение таблицы \ref{analiz:table}}
\centrow Маршрут & \centrow Метод & \centrow Входные данные & \centrow Результат \\ \hline
\finishhead

 \url{/api/notifications/} & GET & -- & Список уведомлений текущего пользователя. \\ \hline
 \url{/api/notifications/{id}/mark_read/} & POST & -- & Пометить уведомление прочитанным. \\ \hline
 \url{/api/notifications/mark_all_read/} & POST & -- & Пометить все уведомления прочитанными. \\ \hline
 \url{ws/notifications/} & WS & JWT в query-параметре. & Push-события уведомлений в реальном времени. \\ \hline
 \url{/api/chats/conversations/} & GET & -- & Список диалогов текущего пользователя. \\ \hline
 \url{/api/chats/conversations/} & POST & participant\_id, item\_id & Создать или получить существующий диалог. \\ \hline
 \url{/api/chats/conversations/{id}/} & GET & -- & Детали диалога с историей сообщений. \\ \hline
 \url{/api/chats/conversations/{id}/messages/} & POST & text, image & Отправить сообщение через HTTP. \\ \hline
 \url{/api/chats/conversations/{id}/mark_read/} & POST & -- & Пометить сообщения прочитанными. \\ \hline
 \url{/api/chats/support-faq/} & GET & -- & Список FAQ для раздела поддержки. \\ \hline
 \url{ws/chat/{chat_id}/} & WS & JWT в query-параметре & Двунаправленный канал чата. \\ \hline
 \url{/api/contracts/by-booking/{booking_id}/} & GET & booking\_id в URL & Договор по бронированию (создать при отсутствии). \\ \hline
 \url{/api/contracts/{id}/sign/} & POST & signer\_name, certificate\_pin & Подписать договор текущим пользователем. \\ \hline
 \url{/api/contracts/} & GET & -- & Список договоров текущего пользователя. \\ \hline
\end{xltabular}

\noindent Пример успешного ответа \verb|GET /api/notifications/|

\begin{verbatim}
[
 {
  "id": 9101,
  "type": "new_message",
  "message": "Пользователь owner отправил вам новое сообщение",
  "metadata": {
   "destination": "chat",
   "chat_id": 15,
   "item_id": 42
  },
  "is_read": false,
  "created_at": "2026-05-15T18:20:10Z",
  "user": 7
 }
]
\end{verbatim}

\noindent Пример успешного ответа \verb|POST /api/notifications/{id}/mark_read/|

\begin{verbatim}
{
 "id": 9101,
 "type": "new_message",
 "message": "Пользователь owner отправил вам новое сообщение",
 "metadata": {
  "destination": "chat",
  "chat_id": 15,
  "item_id": 42
 },
 "is_read": true,
 "created_at": "2026-05-15T18:20:10Z",
 "user": 7
}
\end{verbatim}

\noindent Пример успешного ответа \verb|POST /api/notifications/mark_all_read/|

\begin{verbatim}
{
 "updated": 3
}
\end{verbatim}

\noindent Пример push-события \verb|ws/notifications/|

\begin{verbatim}
{
 "event": "notification_created",
 "notification": {
  "id": 9102,
  "type": "booking_confirmed",
  "message": "Ваше бронирование товара «Дрель» подтверждено",
  "metadata": {
   "destination": "rent_payment",
   "booking_id": 31,
   "item_id": 42
  },
  "is_read": false,
  "created_at": "2026-05-15T18:25:00Z",
  "user": 7
 }
}
\end{verbatim}

\noindent Пример успешного ответа \verb|GET /api/chats/conversations/|

\begin{verbatim}
[
 {
  "id": 15,
  "other_user": {
   "id": 3,
   "username": "owner"
  },
  "item": {
   "id": 42,
   "title": "Дрель"
  },
  "created_at": "2026-05-15T17:10:00Z",
  "unread_count": 2,
  "last_message": {
   "text": "Здравствуйте, товар доступен?",
   "image": null,
   "created_at": "2026-05-15T18:30:00Z",
   "sender": "renter"
  }
 }
]
\end{verbatim}

\noindent Пример успешного ответа \verb|POST /api/chats/conversations/|

\begin{verbatim}
{
 "id": 15,
 "other_user": {
  "id": 3,
  "username": "owner"
 },
 "item": {
  "id": 42,
  "title": "Дрель"
 },
 "created_at": "2026-05-15T17:10:00Z",
 "unread_count": 0,
 "last_message": null
}
\end{verbatim}

\noindent Пример успешного ответа \verb|GET /api/chats/conversations/{id}/|

\begin{verbatim}
{
 "id": 15,
 "other_user": {
  "id": 3,
  "username": "owner"
 },
 "item": {
  "id": 42,
  "title": "Дрель"
 },
 "messages": [
  {
   "id": 120,
   "chat": 15,
   "sender_id": 7,
   "sender_username": "renter",
   "text": "Здравствуйте, товар доступен?",
   "image": null,
   "created_at": "2026-05-15T18:30:00Z",
   "is_read": true
  }
 ]
}
\end{verbatim}

\noindent Пример успешного ответа \verb|POST /api/chats/conversations/{id}/messages/|

\begin{verbatim}
{
 "id": 121,
 "chat": 15,
 "sender_id": 7,
 "sender_username": "renter",
 "text": "Хочу арендовать товар на выходные.",
 "image": null,
 "created_at": "2026-05-15T18:35:00Z",
 "is_read": false
}
\end{verbatim}

\noindent Пример успешного ответа \verb|POST /api/chats/conversations/{id}/mark_read/|

\begin{verbatim}
{
 "status": "ok",
 "message_ids": [120, 121]
}
\end{verbatim}

\noindent Пример успешного ответа \verb|GET /api/chats/support-faq/|

\begin{verbatim}
{
 "questions": [
  {
   "id": "create-listing",
   "question": "Как создать объявление?",
   "answer_title": "Как разместить новое объявление",
   "steps": [
    "Войдите в аккаунт и откройте кнопку \"Объявление\" в верхнем меню.",
    "Укажите название, описание, стоимость аренды за день и выберите категорию.",
    "Добавьте фотографии, чтобы карточка выглядела убедительно для арендаторов.",
    "Нажмите \"Опубликовать\" и дождитесь появления объявления в общем списке."
   ]
  }
 ]
}
\end{verbatim}

\noindent Пример push-события \verb|ws/chat/{chat_id}/|

\begin{verbatim}
{
 "event": "new_message",
 "message": {
  "id": 122,
  "chat": 15,
  "sender_id": 3,
  "sender_username": "owner",
  "text": "Да, товар доступен.",
  "image": null,
  "created_at": "2026-05-15T18:40:00Z",
  "is_read": false
 }
}
\end{verbatim}

\noindent Пример успешного ответа \verb|GET /api/contracts/by-booking/{booking_id}/|

\begin{verbatim}
{
 "id": 501,
 "booking": 31,
 "document_number": "AR-31-A1B2C3D4",
 "file_url": "http://localhost:8000/media/contracts/contract_31_AR-31-A1B2C3D4.pdf",
 "preview_text": "ДОГОВОР АРЕНДЫ № AR-31-A1B2C3D4\n\nПредмет аренды: Дрель\n...",
 "status": "draft",
 "is_signed": false,
 "signed_at": null,
 "created_at": "2026-05-15T18:45:00Z",
 "item_title": "Дрель",
 "start_date": "2026-05-20",
 "end_date": "2026-05-22",
 "total_price": "1500.00",
 "renter_username": "renter",
 "owner_username": "owner",
 "renter_signer_name": "",
 "renter_signature_code": "",
 "renter_signed_at": null,
 "owner_signer_name": "",
 "owner_signature_code": "",
 "owner_signed_at": null,
 "my_role": "renter",
 "my_signed": false
}
\end{verbatim}

\noindent Пример успешного ответа \verb|POST /api/contracts/{id}/sign/|

\begin{verbatim}
{
 "id": 501,
 "booking": 31,
 "document_number": "AR-31-A1B2C3D4",
 "file_url": "http://localhost:8000/media/contracts/contract_31_AR-31-A1B2C3D4.pdf",
 "preview_text": "ДОГОВОР АРЕНДЫ № AR-31-A1B2C3D4\n\nПредмет аренды: Дрель\n...",
 "status": "partially_signed",
 "is_signed": false,
 "signed_at": null,
 "created_at": "2026-05-15T18:45:00Z",
 "item_title": "Дрель",
 "start_date": "2026-05-20",
 "end_date": "2026-05-22",
 "total_price": "1500.00",
 "renter_username": "renter",
 "owner_username": "owner",
 "renter_signer_name": "Иван Иванов",
 "renter_signature_code": "EDS-DEMO-RENTER-31-A1B2C3D4E5",
 "renter_signed_at": "2026-05-15T18:50:00Z",
 "owner_signer_name": "",
 "owner_signature_code": "",
 "owner_signed_at": null,
 "my_role": "renter",
 "my_signed": true
}
\end{verbatim}

\noindent Пример успешного ответа \verb|GET /api/contracts/|

\begin{verbatim}
[
 {
  "id": 501,
  "booking": 31,
  "document_number": "AR-31-A1B2C3D4",
  "file_url": "http://localhost:8000/media/contracts/contract_31_AR-31-A1B2C3D4.pdf",
  "preview_text": "ДОГОВОР АРЕНДЫ № AR-31-A1B2C3D4\n\nПредмет аренды: Дрель\n...",
  "status": "partially_signed",
  "is_signed": false,
  "signed_at": null,
  "created_at": "2026-05-15T18:45:00Z",
  "item_title": "Дрель",
  "start_date": "2026-05-20",
  "end_date": "2026-05-22",
  "total_price": "1500.00",
  "renter_username": "renter",
  "owner_username": "owner",
  "renter_signer_name": "Иван Иванов",
  "renter_signature_code": "EDS-DEMO-RENTER-31-A1B2C3D4E5",
  "renter_signed_at": "2026-05-15T18:50:00Z",
  "owner_signer_name": "",
  "owner_signature_code": "",
  "owner_signed_at": null,
  "my_role": "renter",
  "my_signed": true
 }
]
\end{verbatim}


WebSocket-маршруты объявлены в файлах chats/routing.py и notifications/routing.py и подключаются в asgi.py через URLRouter. WebSocket-маршруты системы представлены в таблице 3.2.
\begin{xltabular}{\textwidth}{|c|X|}
 \caption{WebSocket-маршруты системы \label{websocket:table}}\\ \hline
   Маршрут & \centrow Назначение \\ \hline
 \endfirsthead
 \continuecaption{Продолжение таблицы \ref{websocket:table}}
 ~ Маршрут & \centrow Назначение \\ \hline 
 \finishhead
     /ws/chat/{chat\_id}/ & Обмен сообщениями и событиями прочтения внутри чата. \\ \hline
     /ws/notifications/ & Доставка новых уведомлений текущему пользователю. \\ \hline
\end{xltabular}

\noindent Пример отправки сообщения через \verb|/ws/chat/{chat_id}/|

\begin{verbatim}
{
 "message": "Здравствуйте, товар еще доступен?"
}
\end{verbatim}

\noindent Пример события нового сообщения из \verb|/ws/chat/{chat_id}/|

\begin{verbatim}
{
 "event": "new_message",
 "message": {
  "id": 122,
  "chat": 15,
  "sender_id": 7,
  "sender_username": "renter",
  "text": "Здравствуйте, товар еще доступен?",
  "image": null,
  "created_at": "2026-05-15T18:40:00Z",
  "is_read": false
 }
}
\end{verbatim}

\noindent Пример отправки события прочтения через \verb|/ws/chat/{chat_id}/|

\begin{verbatim}
{
 "action": "mark_read"
}
\end{verbatim}

\noindent Пример события о прочтении сообщений из \verb|/ws/chat/{chat_id}/|

\begin{verbatim}
{
 "event": "messages_read",
 "message_ids": [120, 121],
 "read_by_user_id": 7
}
\end{verbatim}

\noindent Пример события нового уведомления из \verb|/ws/notifications/|

\begin{verbatim}
{
 "event": "notification_created",
 "notification": {
  "id": 9102,
  "type": "booking_confirmed",
  "message": "Ваше бронирование товара «Дрель» подтверждено",
  "metadata": {
   "destination": "rent_payment",
   "booking_id": 31,
   "item_id": 42
  },
  "is_read": false,
  "created_at": "2026-05-15T18:25:00Z",
  "user": 7
 }
}
\end{verbatim}


\subsubsection{Детализация функциональной подсистемы уведомлений}
Подсистема уведомлений предназначена для информирования пользователей о важных событиях. В проекте выделены следующие типы уведомлений:
\begin{itemize}
    \item новая заявка на аренду;
    \item подтверждение бронирования;
    \item отмена бронирования;
    \item подтверждение оплаты;
    \item напоминание о возврате;
    \item новое сообщение;
    \item новый отзыв.
\end{itemize}

Модель уведомления содержит пользователя-получателя, тип уведомления, текст, дополнительные метаданные, признак прочтения и дату создания. Метаданные используются для передачи служебной информации: идентификатора бронирования, идентификатора товара, идентификатора чата и направления перехода в интерфейсе.

Создание уведомлений вынесено в сервисный слой. Это позволяет вызывать уведомления из разных частей проекта без дублирования логики. Например, при создании бронирования вызывается функция уведомления арендодателя, при подтверждении бронирования - уведомление арендатора, а при отправке сообщения - уведомление всех участников чата, кроме отправителя.

Важной особенностью является то, что сохранение уведомления в базе данных и realtime-доставки разделены. Если WebSocket-канал недоступен, уведомление всё равно сохраняется. Это повышает надёжность системы: пользователь сможет увидеть уведомление позже через REST API.

Email-отправка также реализована как дополнительный канал. Ошибка отправки email не должна нарушать основной бизнес-процесс. Поэтому обработка исключений в почтовом слое не приводит к падению сервера. Для SMS-канала может использоваться аналогичный подход: отдельная функция отправки через провайдера, вызываемая из того же сервисного слоя.

\subsubsection{Детализация функциональной подсистемы чатов}
Подсистема чатов предназначена для общения арендатора и арендодателя. Диалог может быть связан с конкретным предметом аренды. Это важно для пользовательского опыта, так как у одного и того же арендодателя может быть несколько товаров, и переписка по каждому предмету должна быть отделена.

При создании диалога сервер проверяет:
\begin{enumerate}
    \item Указан ли участник переписки.
    \item Существует ли такой пользователь.
    \item Не пытается ли пользователь создать чат с самим собой.
    \item Существует ли указанный товар.
    \item Принадлежит ли товар выбранному участнику.
    \item Не существует ли уже диалог между этими пользователями по этому товару.
\end{enumerate}

Если подходящий диалог уже есть, возвращается существующая запись. Если нет, создаётся новый чат, и оба пользователя добавляются в связь многие-ко-многим.

Сообщение может содержать текст или изображение. Сериализатор проверяет, что хотя бы одно из этих полей заполнено. Пустые сообщения отклоняются. После сохранения сообщения система создаёт уведомление получателю и передаёт событие в WebSocket-группу соответствующего чата.

Для чтения сообщений реализована отдельные алгоритмы. Когда пользователь открывает диалог или отправляет запрос на отметку прочтения, сервер находит все непрочитанные сообщения, отправленные другим участником, отмечает их как прочитанные и рассылает событие messages\_read в WebSocket-группу. Это позволяет синхронизировать состояние переписки между участниками.

\subsubsection{Детализация функциональной подсистемы договоров}
Подсистема договоров предназначена для автоматического формирования документа аренды. Договор связан с бронированием отношением один-к-одному. Это означает, что для одного бронирования может существовать только один договор. Такой подход исключает появление нескольких конкурирующих документов по одной сделке.

Договор создаётся только для бронирований со статусом confirmed или completed. Для заявок в статусе ожидания договор недоступен, так как сделка ещё не согласована сторонами. При обращении к API по бронированию сервер проверяет, является ли текущий пользователь участником сделки: арендатором или владельцем предмета. Если пользователь не относится к сделке, договор ему не выдаётся.

Номер договора формируется автоматически и содержит префикс проекта, идентификатор бронирования и случайный уникальный фрагмент. Такой номер удобен для визуальной идентификации документа и снижает риск совпадений.

Текст договора формируется из данных системы: предмет аренды, стороны договора, даты аренды, сумма, статус бронирования и сведения о подписании. После формирования текст преобразуется в PDF-файл. При подписании одной из сторон в договор добавляются имя подписанта, дата подписания и демонстрационный код подписи. После подписи обеих сторон договор получает статус подписанного.

\subsubsection{Детализация операций доступа к данным}
В подсистеме уведомлений основным запросом является получение уведомлений текущего пользователя. Запрос фильтруется по user\_id, что исключает доступ к чужим уведомлениям. Сортировка выполняется по дате создания в обратном порядке, чтобы новые уведомления отображались первыми.

Операции изменения ограничены отметкой уведомлений как прочитанных. Пользователь может изменить только свои уведомления, так как выборка viewset изначально фильтруется по текущему пользователю.

В подсистеме чатов выборка диалогов также ограничивается текущим пользователем. Список диалогов сортируется по дате последнего сообщения. Для этого используется агрегирование по сообщениям, что позволяет показывать сначала наиболее активные переписки.

При получении конкретного диалога сервер проверяет, что текущий пользователь входит в список участников. Если пользователь не является участником, возвращается ошибка. Такая проверка применяется и при отправке сообщения.

В подсистеме договоров выборка ограничивается договорами, где текущий пользователь является арендатором или владельцем предмета. Это реализуется через условие по бронированию и владельцу товара. При получении договора по бронированию дополнительно проверяется статус бронирования. Договор не выдаётся для неподтверждённых заявок.

\subsubsection{Нормализация}
Схема данных соответствует требованиям нормализации. Повторяющиеся данные вынесены в отдельные сущности. Например, сведения о пользователях не дублируются в уведомлениях, чатах или договорах, а подключаются через внешние ключи. Данные о товаре также не копируются в чат, а связываются через внешний ключ.

Договор хранит только те данные, которые относятся к состоянию самого договора: номер, файл, подписи, даты подписания. Основная информация о сделке берётся из бронирования. Это уменьшает риск противоречий между договором и бронированием.

В уведомлениях используется поле metadata, так как оно содержит вспомогательные данные для навигации и не является отдельной нормализованной сущностью. Такой подход оправдан, поскольку метаданные зависят от типа уведомления и не участвуют в сложных реляционных связях.

\subsection{Надёжность, производительность и масштабирование}

\subsubsection{Надёжность операций уведомлений}
Надёжность подсистемы уведомлений обеспечивается разделением критических и некритических операций. Критической операцией является создание записи уведомления в базе данных. Некритическими являются realtime-доставка и внешняя отправка email. Если внешний канал временно недоступен, основное уведомление не теряется.

Такой подход важен для пользовательских событий. Например, если при создании бронирования почтовый сервер недоступен, арендодатель всё равно получит уведомление в системе после входа в аккаунт. Это снижает зависимость бизнес-логики от внешних сервисов.

\subsubsection{Надёжность чатов}
Для чатов надёжность обеспечивается сохранением сообщения в базе данных до рассылки WebSocket-события. Если пользователь-получатель не находится онлайн, сообщение не теряется и будет доступно при следующем открытии диалога.

Проверка членства пользователя в чате выполняется как при WebSocket-подключении, так и при HTTP-операциях. Это защищает переписку от несанкционированного доступа.

\subsubsection{Надёжность формирования договоров}
Договор формируется на основании уже сохранённого бронирования. Это исключает ситуацию, когда документ создаётся без связи с реальной сделкой. Ограничение один-к-одному между бронированием и договором предотвращает появление нескольких договоров по одной заявке.

При подписании договора проверяется роль пользователя. Арендатор не может подписаться за арендодателя, а сторонний пользователь не может подписать договор вообще. После каждой подписи PDF-файл пересоздаётся, чтобы его содержимое соответствовало текущему состоянию модели.

\subsubsection{Потенциал масштабирования}
Текущая архитектура допускает дальнейшее развитие:
\begin{itemize}
    \item выделение сервиса уведомлений в отдельный микросервис;
    \item подключение очереди задач для фоновой email-рассылки;
    \item перенос генерации PDF в отдельный воркер;
    \item постраничная загрузка сообщений в чатах;
    \item хранение файлов договоров во внешнем объектном хранилище;
    \item масштабирование WebSocket-сервера с использованием Redis;
\end{itemize}

\subsection{Информационная безопасность технического проекта}
\subsubsection{Профиль угроз}
Для разработанных модулей наиболее актуальны следующие угрозы:
\begin{itemize}
    \item доступ пользователя к чужим уведомлениям;
    \item доступ к чужой переписке;
    \item подключение к WebSocket чужого чата;
    \item подписание договора пользователем, не являющимся участником сделки;
    \item отправка пустых или некорректных сообщений;
    \item раскрытие персональных данных участников аренды.
\end{itemize}

Так как система работает с пользовательскими сообщениями, email-адресами, договорами и данными сделок, контроль доступа является одним из ключевых требований.

\subsubsection{Меры защиты}
В проекте применяются следующие меры защиты:
\begin{enumerate}
    \item Авторизация через JWT. Доступ к основным API предоставляется только аутентифицированным пользователям.
    \item Фильтрация данных по текущему пользователю. Уведомления, чаты и договоры выбираются не из общего списка, а через условия, связанные с пользователем запроса.
    \item Проверка членства в чате. Пользователь может открыть WebSocket-соединение только с тем чатом, где он является участником.
    \item Проверка участия в сделке. Договор доступен только арендатору и арендодателю.
    \item Валидация сообщений. Пустые сообщения без текста и изображения отклоняются.
    \item Ограничение статусов договора. Договор формируется только для подтверждённых или завершённых бронирований.
    \item Разделение основной операции и внешней доставки уведомлений. Ошибка почтового сервера не раскрывается пользователю и не нарушает бизнес-процесс.
    \item Использование ORM. Работа с базой данных выполняется через Django ORM, что снижает риск SQL-инъекций.
\end{enumerate}
